% SPDX-License-Identifier: CC-BY-SA-4.0
% Author: Matthieu Perrin

\begingroup

\begin{exercice}[Propriétés des automates]\label{exo:lexing/determinisation/determinisation}

  \begin{question}
  \item
    Donner l'automate déterministe équivalent à l'automate suivant
    par la méthode de construction des sous-ensembles de Rabin et Scott.

    \begin{tikzpicture}[automaton, baseline=(1)]
      \state[initial]   (1)  at (0,2) {$1$}; 
      \state            (2)  at (1,2) {$2$}; 
      \state            (3)  at (0,1) {$3$}; 
      \state            (4)  at (2,2) {$4$}; 
      \state            (5)  at (3,2) {$5$}; 
      \state            (6)  at (4,2) {$6$}; 
      \state            (7)  at (1,1) {$7$}; 
      \state            (8)  at (2,1) {$8$}; 
      \state[initial]   (9)  at (0,0) {$9$}; 
      \state            (10) at (1,0) {$10$}; 
      \state[accepting] (11) at (4,1) {$11$}; 

      \path (1)  edge                node[swap] {$\varepsilon$} (2);
      \path (2)  edge                node[swap] {$\varepsilon$} (4);
      \path (4)  edge                node       {$a$}           (5);
      \path (5)  edge                node[swap] {$\varepsilon$} (6);
      \path (6)  edge                node       {$\varepsilon$} (11);
      \path (1)  edge                node       {$\varepsilon$} (3);
      \path (3)  edge                node       {$a$}           (7);
      \path (7)  edge                node       {$b$}           (8);
      \path (8)  edge                node       {$\varepsilon$} (11);
      \path (9)  edge                node       {$\varepsilon$} (10);
      \path (10) edge                node[swap] {$b$}           (8);
      \path (2)  edge[bend left=8mm] node       {$\varepsilon$} (6);
      \path (5)  edge[bend left=8mm] node       {$\varepsilon$} (4);
    \end{tikzpicture}

  \end{question}
  
\end{exercice}

\begin{correction}

  \textbf{Rappels de vocabulaire.}
  \begin{description}
  \item[Unitaire] : un seul état initial.
  \item[Standard] : unitaire sans transition entrante dans l'état initial.
  \item[Normalisé] : standard + un seul état final sans transition sortante.
  \item[$\varepsilon$-libre] : aucune $\varepsilon$-transition.
  \item[Déterministe] : Unitaire, $\varepsilon$-libre, et au plus une transition par lettre et par état.
  \item[Complet] : au moins une transition par lettre et par état.
  \end{description}
  
  \begin{question}
  \item 

    \underline{Propriétés.}
    \begin{itemize}
    \item Unitaire : \textbf{non} (états finaux $\{1,6\}$).
    \item Standard : \textbf{oui} (un seul initial $1$, aucune entrée vers $1$).
    \item Normalisé : \textbf{non} (plusieurs finaux et transitions sortantes depuis $1$).
    \item $\varepsilon$-libre : \textbf{oui}.
    \item Déterministe : \textbf{non} (depuis $6$ deux transitions en $0$; depuis $1$ deux sorties en $0$ : boucle et $1\to4$).
    \item Complet : \textbf{non} (p.ex. $3$ n’a pas de $1$; $5$ n’a pas de $0$; $6$ n’a pas de $1$).
    \end{itemize}

    \underline{Complétion.}
    Ajouter un puits $\bot$ bouclé sur $0,1$, et compléter :
    \[
    3 \xrightarrow{1} \bot,\quad
    5 \xrightarrow{0} \bot,\quad
    6 \xrightarrow{1} \bot.
    \]
    (Tous les autres états ont déjà au moins une sortie sur chaque lettre.)

    \underline{Déterminisation (méthode des sous-ensembles).}
    États AFD = sous-ensembles d’états; initial $\{1\}$; finaux = sous-ensembles contenant $1$ ou $6$.
    Table (transitions $0/1$) — on ne liste que les ensembles atteignables :
    \[
    \renewcommand{\arraystretch}{1.2}
    \begin{array}{c|c|c}
      X & \delta(X,0) & \delta(X,1)\\\hline
      \{1\}     & \{1,4\} & \{2\}\\
      \{1,4\}   & \{1,4\} & \{2,4\}\\
      \{2\}     & \{3\}   & \{4\}\\
      \{2,4\}   & \{3,4\} & \{4\}\\
      \{3\}     & \{6\}   & \emptyset\\
      \{4\}     & \{4\}   & \{4\}\\
      \{3,4\}   & \{4,6\} & \{4\}\\
      \{6\}     & \{3,5\} & \emptyset\\
      \{4,6\}   & \{3,4,5\} & \{4\}\\
      \{3,5\}   & \{6\}   & \{5\}\\
      \{5\}     & \emptyset & \{5\}\\
      \{3,4,5\} & \{4,6\} & \{4,5\}\\
      \{4,5\}   & \{4\}   & \{4,5\}
    \end{array}
    \]
    (On ajoute l’état puits $\emptyset$ si l’on souhaite un DFA complet.)
    États \textit{acceptants} : ceux qui contiennent $1$ ou $6$ (ex. $\{1\},\{1,4\},\{6\},\{4,6\},\{3,5\}$…).

    \medskip

    % ===========================
    % (2)
    % ===========================
  \item \textbf{Automate (2) — alphabet $\{a,b\}$.}

    \underline{Propriétés.}
    \begin{itemize}
    \item Unitaire : \textbf{oui} (un seul final $4$).
    \item Standard : \textbf{non} (des entrées vers l’état initial $1$ : $2\xrightarrow{a}1$, $3\xrightarrow{b}1$).
    \item Normalisé : \textbf{non} (l’état final $4$ a des sorties).
    \item $\varepsilon$-libre : \textbf{oui}.
    \item Déterministe : \textbf{oui}.
    \item Complet : \textbf{oui} (chaque état a une sortie en $a$ et en $b$).
    \end{itemize}

    \underline{Complétion.} Aucune (déjà complet).
    \underline{Déterminisation.} Inutile (déjà DFA).

    \medskip

    % ===========================
    % (3)
    % ===========================
  \item \textbf{Automate (3) — alphabet $\{a,b\}$.}

    \underline{Propriétés.}
    \begin{itemize}
    \item Unitaire : \textbf{oui} (final $\{2\}$).
    \item Standard : \textbf{non} (boucle sur l’état initial $0$).
    \item Normalisé : \textbf{non} (final $2$ a des sorties).
    \item $\varepsilon$-libre : \textbf{oui}.
    \item Déterministe : \textbf{non} (depuis $0$: deux sorties en $a$ — boucle et $0\to1$; idem $1$ en $a$).
    \item Complet : \textbf{oui} (chaque état a au moins une sortie sur $a$ et sur $b$).
    \end{itemize}

    \underline{Complétion.} Aucune (déjà complet).

    \underline{Déterminisation (sous-ensembles).}
    Initial $\{0\}$; finaux = sous-ensembles contenant $2$.
    \[
    \renewcommand{\arraystretch}{1.2}
    \begin{array}{c|c|c}
      X & \delta(X,a) & \delta(X,b)\\\hline
      \{0\}     & \{0,1\}   & \{0\}\\
      \{0,1\}   & \{0,1,2\} & \{0,1\}\\
      \{0,1,2\} & \{0,1,2\} & \{0,1,2\}
    \end{array}
    \]
    DFA obtenu (équivalent) : états $\{0\}$, $\{0,1\}$, $\{0,1,2\}$; seul $\{0,1,2\}$ est acceptant.

    \medskip

    % ===========================
    % (4)
    % ===========================
  \item \textbf{Automate (4) — alphabet $\{a,b\}$.}

    \underline{Propriétés.}
    \begin{itemize}
    \item Unitaire : \textbf{oui} (final $2$).
    \item Standard : \textbf{oui} (un seul initial $0$, pas d’entrée vers $0$).
    \item Normalisé : \textbf{non} (le final $2$ a une boucle sortante).
    \item $\varepsilon$-libre : \textbf{oui}.
    \item Déterministe : \textbf{oui}.
    \item Complet : \textbf{non} ($0$ n’a pas de $b$, $1$ n’a pas de $a$, $2$ n’a pas de $b$).
    \end{itemize}

    \underline{Complétion.}
    Ajouter un puits $\bot$ bouclé en $a,b$, et compléter :
    \[
    0 \xrightarrow{b} \bot,\qquad
    1 \xrightarrow{a} \bot,\qquad
    2 \xrightarrow{b} \bot.
    \]
    \underline{Déterminisation.} Inutile (déjà DFA).

  \end{question}

\end{correction}



\begin{correction}

  \begin{question}
    
  \item Correction détaillée pour l'automate A.
    
    \textbf{Phase 1.}
    La construction de l'AFD équivalent commence par la détermination de l'état initial.
    Celui-ci est construit à partir de l' $\varepsilon$-fermeture de tous les états initiaux.
    Donc ici, $I'=\{0\}$. A partir de là, on construit la table des transitions.
    
    $$\begin{array}{| l | c | c|}
      \hline
      \  \mu & a & b \\
      \hline
      \  A = \{0\} = I' &  &  \\
      \hline
    \end{array}$$

    \textbf{Phase 2.}
    On regarde vers quels états l'automate se trouve pour chacun des symboles de l'alphabet à partir de $I'$.
    Normalement, il faut faire l'$\varepsilon$-fermeture sur les états obtenus mais dans cet automate il n'y a aucune $\varepsilon$-transition.
    On les ajoute à la table s'ils n'y sont pas déjà.

    $$\begin{array}{| l | c | c|}
      \hline
      \  \mu & a & b \\
      \hline
      \  A = \{0\} = I' & \{0,1\} & \{0\} \\
      \  B = \{0,1\} &  &  \\
      \hline
    \end{array}$$


    \textbf{Phase 3.}
    Puis on fait la même chose pour chacun des états pas encore traités.
    $$\begin{array}{| l | c | c|}
      \hline
      \  \mu & a & b \\
      \hline
      \  A = \{0\} = I' & \{0,1\} & \{0\} \\
      \  B = \{0,1\} &  \{0,1,2 \}& \{0,1\}  \\
      \  C = \{0,1,2\} &  &  \\
      \hline
    \end{array}$$

    \textbf{Phase 4.}
    C contient l'état 2 qui est final dans l'automate de départ. Donc C est un état final dans l'AFD obtenu.

    $$\begin{array}{| l | c | c|}
      \hline
      \  \mu & a & b \\
      \hline
      \  A = \{0\} = I' & \{0,1\} & \{0\} \\
      \  B = \{0,1\} &  \{0,1,2 \} & \{0,1\}  \\
      \  C = \{0,1,2\} \in F &  \{0,1,2 \} &  \{0,1,2 \} \\
      \hline
    \end{array}$$

    Finalement, le tableau est rempli, l'automate est terminé. On obtient l'automate suivant :

    \begin{center}
      \begin{tikzpicture}[automaton]
        \state[initial  ] (0) at (0,0) {$A$}; 
        \state            (1) at (1,0) {$B$}; 
        \state[accepting] (2) at (2,0) {$C$}; 

        \path (0) edge[loop above] node {$b$}    (0);
        \path (1) edge[loop above] node {$b$}    (1);
        \path (2) edge[loop above] node {$a, b$} (2);
        \path (0) edge             node {$a$}    (1);
        \path (1) edge             node {$a$}    (2);
      \end{tikzpicture}
    \end{center}

  \item Automate $B$ déterminisé :\\
    \begin{minipage}[t]{.4\textwidth}
      \begin{itemize}
      \item $A = \{1, 2, 3, 4, 6, 9, 10, 11\}$
      \item $B = \{8, 11\}$
      \item $C = \{4, 5, 6, 7, 11\}$
      \item $D = \{4, 5, 6, 11\}$
      \end{itemize}
    \end{minipage}%
    \begin{minipage}[t]{.3\textwidth}
      \begin{tikzpicture}[automaton]
        \state[accepting, initial] (A) at (0,1) {$A$}; 
        \state[accepting]          (B) at (1,1) {$B$}; 
        \state[accepting]          (C) at (0,0) {$C$}; 
        \state[accepting]          (D) at (1,0) {$D$}; 

        \path (D) edge[loop right] node {$a$}  (D);
        \path (A) edge             node {$b$} (B);
        \path (A) edge             node {$a$} (C);
        \path (C) edge             node {$b$} (B);
        \path (C) edge             node {$a$} (D);
      \end{tikzpicture}
    \end{minipage}

  \end{question}

\end{correction}

\begin{correction}

  \begin{question}
  \item {\small
    \begin{tabular}[t]{| l | c c c c |}
      \hline
      & A & B & C  & D \\
      \hline
      Accessible          & $\{1,2,3,4\}$ & $\{1, 2, 3, 4, 5, 7\}$  & $\{1, 2, 3, 4, 5, 6\}$  & $\{1, 2, 3, 4, 5, 6, 7, 8, 9, 10, 11\}$ \\
      Inaccessible        & $\emptyset$   & $\{6\}$                 & $\emptyset$             & $\emptyset$                             \\
      Co-accessible       & $\{1,2,3,4\}$ & $\{1, 2, 5\}$           & $\{1, 2, 3, 6\}$        & $\{1, 2, 3, 4, 5, 6, 7, 8, 9, 10, 11\}$ \\
      Stérile             & $\emptyset$   & $\{3, 4, 6, 7\}$        & $\{4, 5\}$              & $\emptyset$                             \\
      Utile               & $\{1,2,3,4\}$ & $\{1, 2, 5\}$           & $\{1, 2, 3, 6\}$        & $\{1, 2, 3, 4, 5, 6, 7, 8, 9, 10, 11\}$ \\
      Inutile             & $\emptyset$   & $\{3, 4, 6, 7\}$        & $\{4, 5\}$              & $\emptyset$                             \\
      Complet             & oui           & non                     & non                     & non                                     \\
      Emondé              & oui           & non                     & non                     & oui                                     \\
      Unitaire            & oui           & oui                     & oui                     & oui                                     \\
      Standard            & non           & oui                     & non                     & oui                                     \\
      Normalisé           & non           & non                     & non                     & non                                     \\
      $\varepsilon$-libre & oui           & oui                     & oui                     & non                                     \\
      Déterministe        & oui           & non                     & non                     & non                                     \\
      \hline
  \end{tabular}}
  \item 
    \begin{minipage}[t]{.4\textwidth}
      $B$ sur l'alphabet $\{a,b,c,d\}$ :\\

      \scalebox{.9}{\begin{tikzpicture}[shorten >=1pt,node distance=1.7cm,on grid,auto]
          \node[state, initial, initial text=, accepting] (1)              {$1$}; 
          \node[state, accepting]                         (2) [right=of 1] {$2$}; 
          \node[state, accepting]                         (5) [below=of 2] {$5$}; 

          \path[-latex] (2) edge[loop above, looseness=5] node {$a$}    (2);
          \path[-latex] (5) edge[loop below, looseness=5] node {$b, c$}    (5);
          \path[-latex] (1) edge node {$a$} (2);
          \path[-latex] (1) edge node[swap] {$b, c$} (5);
          \path[-latex] (2) edge[bend left=6mm] node {$b$} (5);
          \path[-latex] (5) edge[bend left=6mm] node {$a$} (2);
      \end{tikzpicture}}
    \end{minipage}%
    \begin{minipage}[t]{.4\textwidth}
      $C$ sur l'alphabet $\{0,1,2,3,4,5,6,7,8,9\}$ :\\
      \scalebox{.9}{\begin{tikzpicture}[shorten >=1pt,node distance=1.7cm,on grid,auto]
          \node[state, initial, initial text=, accepting] (1)              {$1$}; 
          \node[state]                                    (2) [right=of 1] {$2$}; 
          \node[state]                                    (3) [right=of 2] {$3$}; 
          \node[state, accepting]                         (6) [below=of 3] {$6$}; 

          \path[-latex] (1) edge[loop above, looseness=5] node {$0$}    (1);
          \path[-latex] (1) edge node {$1$} (2);
          \path[-latex] (2) edge node {$0$} (3);
          \path[-latex] (3) edge[bend left=6mm] node {$0$} (6);
          \path[-latex] (6) edge[bend left=6mm] node {$0$} (3);
      \end{tikzpicture}}
    \end{minipage}
  \end{question}

\end{correction}

\endgroup
\endinput

% 
% 
% 
% 
% 
% 
% 
% 
% 
% 
% 
% 
% 
% 
% 
% 
% 
% 
%\begin{correction}
% 
%  \textbf{Rappels de vocabulaire.}
%  \begin{description}
%  \item[Unitaire] : un seul état initial.
%  \item[Standard] : unitaire sans transition entrante dans l'état initial.
%  \item[Normalisé] : standard + un seul état final sans transition sortante.
%  \item[$\varepsilon$-libre] : aucune $\varepsilon$-transition.
%  \item[Déterministe] : Unitaire, $\varepsilon$-libre, et au plus une transition par lettre et par état.
%  \item[Complet] : au moins une transition par lettre et par état.
%  \end{description}
%  
%  \begin{question}
%  \item 
% 
%    \underline{Propriétés.}
%    \begin{itemize}
%    \item Unitaire : \textbf{non} (états finaux $\{1,6\}$).
%    \item Standard : \textbf{oui} (un seul initial $1$, aucune entrée vers $1$).
%    \item Normalisé : \textbf{non} (plusieurs finaux et transitions sortantes depuis $1$).
%    \item $\varepsilon$-libre : \textbf{oui}.
%    \item Déterministe : \textbf{non} (depuis $6$ deux transitions en $0$; depuis $1$ deux sorties en $0$ : boucle et $1\to4$).
%    \item Complet : \textbf{non} (p.ex. $3$ n’a pas de $1$; $5$ n’a pas de $0$; $6$ n’a pas de $1$).
%    \end{itemize}
% 
%    \underline{Complétion.}
%    Ajouter un puits $\bot$ bouclé sur $0,1$, et compléter :
%    \[
%    3 \xrightarrow{1} \bot,\quad
%    5 \xrightarrow{0} \bot,\quad
%    6 \xrightarrow{1} \bot.
%    \]
%    (Tous les autres états ont déjà au moins une sortie sur chaque lettre.)
% 
%    \underline{Déterminisation (méthode des sous-ensembles).}
%    États AFD = sous-ensembles d’états; initial $\{1\}$; finaux = sous-ensembles contenant $1$ ou $6$.
%    Table (transitions $0/1$) — on ne liste que les ensembles atteignables :
%    \[
%    \renewcommand{\arraystretch}{1.2}
%    \begin{array}{c|c|c}
%      X & \delta(X,0) & \delta(X,1)\\\hline
%      \{1\}     & \{1,4\} & \{2\}\\
%      \{1,4\}   & \{1,4\} & \{2,4\}\\
%      \{2\}     & \{3\}   & \{4\}\\
%      \{2,4\}   & \{3,4\} & \{4\}\\
%      \{3\}     & \{6\}   & \emptyset\\
%      \{4\}     & \{4\}   & \{4\}\\
%      \{3,4\}   & \{4,6\} & \{4\}\\
%      \{6\}     & \{3,5\} & \emptyset\\
%      \{4,6\}   & \{3,4,5\} & \{4\}\\
%      \{3,5\}   & \{6\}   & \{5\}\\
%      \{5\}     & \emptyset & \{5\}\\
%      \{3,4,5\} & \{4,6\} & \{4,5\}\\
%      \{4,5\}   & \{4\}   & \{4,5\}
%    \end{array}
%    \]
%    (On ajoute l’état puits $\emptyset$ si l’on souhaite un DFA complet.)
%    États \textit{acceptants} : ceux qui contiennent $1$ ou $6$ (ex. $\{1\},\{1,4\},\{6\},\{4,6\},\{3,5\}$…).
% 
%    \medskip
% 
%    % ===========================
%    % (2)
%    % ===========================
%  \item \textbf{Automate (2) — alphabet $\{a,b\}$.}
% 
%    \underline{Propriétés.}
%    \begin{itemize}
%    \item Unitaire : \textbf{oui} (un seul final $4$).
%    \item Standard : \textbf{non} (des entrées vers l’état initial $1$ : $2\xrightarrow{a}1$, $3\xrightarrow{b}1$).
%    \item Normalisé : \textbf{non} (l’état final $4$ a des sorties).
%    \item $\varepsilon$-libre : \textbf{oui}.
%    \item Déterministe : \textbf{oui}.
%    \item Complet : \textbf{oui} (chaque état a une sortie en $a$ et en $b$).
%    \end{itemize}
% 
%    \underline{Complétion.} Aucune (déjà complet).
%    \underline{Déterminisation.} Inutile (déjà DFA).
% 
%    \medskip
% 
%    % ===========================
%    % (3)
%    % ===========================
%  \item \textbf{Automate (3) — alphabet $\{a,b\}$.}
% 
%    \underline{Propriétés.}
%    \begin{itemize}
%    \item Unitaire : \textbf{oui} (final $\{2\}$).
%    \item Standard : \textbf{non} (boucle sur l’état initial $0$).
%    \item Normalisé : \textbf{non} (final $2$ a des sorties).
%    \item $\varepsilon$-libre : \textbf{oui}.
%    \item Déterministe : \textbf{non} (depuis $0$: deux sorties en $a$ — boucle et $0\to1$; idem $1$ en $a$).
%    \item Complet : \textbf{oui} (chaque état a au moins une sortie sur $a$ et sur $b$).
%    \end{itemize}
% 
%    \underline{Complétion.} Aucune (déjà complet).
% 
%    \underline{Déterminisation (sous-ensembles).}
%    Initial $\{0\}$; finaux = sous-ensembles contenant $2$.
%    \[
%    \renewcommand{\arraystretch}{1.2}
%    \begin{array}{c|c|c}
%      X & \delta(X,a) & \delta(X,b)\\\hline
%      \{0\}     & \{0,1\}   & \{0\}\\
%      \{0,1\}   & \{0,1,2\} & \{0,1\}\\
%      \{0,1,2\} & \{0,1,2\} & \{0,1,2\}
%    \end{array}
%    \]
%    DFA obtenu (équivalent) : états $\{0\}$, $\{0,1\}$, $\{0,1,2\}$; seul $\{0,1,2\}$ est acceptant.
% 
%    \medskip
% 
%    % ===========================
%    % (4)
%    % ===========================
%  \item \textbf{Automate (4) — alphabet $\{a,b\}$.}
% 
%    \underline{Propriétés.}
%    \begin{itemize}
%    \item Unitaire : \textbf{oui} (final $2$).
%    \item Standard : \textbf{oui} (un seul initial $0$, pas d’entrée vers $0$).
%    \item Normalisé : \textbf{non} (le final $2$ a une boucle sortante).
%    \item $\varepsilon$-libre : \textbf{oui}.
%    \item Déterministe : \textbf{oui}.
%    \item Complet : \textbf{non} ($0$ n’a pas de $b$, $1$ n’a pas de $a$, $2$ n’a pas de $b$).
%    \end{itemize}
% 
%    \underline{Complétion.}
%    Ajouter un puits $\bot$ bouclé en $a,b$, et compléter :
%    \[
%    0 \xrightarrow{b} \bot,\qquad
%    1 \xrightarrow{a} \bot,\qquad
%    2 \xrightarrow{b} \bot.
%    \]
%    \underline{Déterminisation.} Inutile (déjà DFA).
% 
%  \end{question}
% 
%\end{correction}
% 
% 
% 
%\begin{correction}
% 
%  \begin{question}
%    
%  \item Correction détaillée pour l'automate A.
%    
%    \textbf{Phase 1.}
%    La construction de l'AFD équivalent commence par la détermination de l'état initial.
%    Celui-ci est construit à partir de l' $\varepsilon$-fermeture de tous les états initiaux.
%    Donc ici, $I'=\{0\}$. A partir de là, on construit la table des transitions.
%    
%    $$\begin{array}{| l | c | c|}
%      \hline
%      \  \mu & a & b \\
%      \hline
%      \  A = \{0\} = I' &  &  \\
%      \hline
%    \end{array}$$
% 
%    \textbf{Phase 2.}
%    On regarde vers quels états l'automate se trouve pour chacun des symboles de l'alphabet à partir de $I'$.
%    Normalement, il faut faire l'$\varepsilon$-fermeture sur les états obtenus mais dans cet automate il n'y a aucune $\varepsilon$-transition.
%    On les ajoute à la table s'ils n'y sont pas déjà.
% 
%    $$\begin{array}{| l | c | c|}
%      \hline
%      \  \mu & a & b \\
%      \hline
%      \  A = \{0\} = I' & \{0,1\} & \{0\} \\
%      \  B = \{0,1\} &  &  \\
%      \hline
%    \end{array}$$
% 
% 
%    \textbf{Phase 3.}
%    Puis on fait la même chose pour chacun des états pas encore traités.
%    $$\begin{array}{| l | c | c|}
%      \hline
%      \  \mu & a & b \\
%      \hline
%      \  A = \{0\} = I' & \{0,1\} & \{0\} \\
%      \  B = \{0,1\} &  \{0,1,2 \}& \{0,1\}  \\
%      \  C = \{0,1,2\} &  &  \\
%      \hline
%    \end{array}$$
% 
%    \textbf{Phase 4.}
%    C contient l'état 2 qui est final dans l'automate de départ. Donc C est un état final dans l'AFD obtenu.
% 
%    $$\begin{array}{| l | c | c|}
%      \hline
%      \  \mu & a & b \\
%      \hline
%      \  A = \{0\} = I' & \{0,1\} & \{0\} \\
%      \  B = \{0,1\} &  \{0,1,2 \} & \{0,1\}  \\
%      \  C = \{0,1,2\} \in F &  \{0,1,2 \} &  \{0,1,2 \} \\
%      \hline
%    \end{array}$$
% 
%    Finalement, le tableau est rempli, l'automate est terminé. On obtient l'automate suivant :
% 
%    \begin{center}
%      \begin{tikzpicture}[automaton]
%        \state[initial  ] (0) at (0,0) {$A$}; 
%        \state            (1) at (1,0) {$B$}; 
%        \state[accepting] (2) at (2,0) {$C$}; 
% 
%        \path (0) edge[loop above] node {$b$}    (0);
%        \path (1) edge[loop above] node {$b$}    (1);
%        \path (2) edge[loop above] node {$a, b$} (2);
%        \path (0) edge             node {$a$}    (1);
%        \path (1) edge             node {$a$}    (2);
%      \end{tikzpicture}
%    \end{center}
% 
%  \item Automate $B$ déterminisé :\\
%    \begin{minipage}[t]{.4\textwidth}
%      \begin{itemize}
%      \item $A = \{1, 2, 3, 4, 6, 9, 10, 11\}$
%      \item $B = \{8, 11\}$
%      \item $C = \{4, 5, 6, 7, 11\}$
%      \item $D = \{4, 5, 6, 11\}$
%      \end{itemize}
%    \end{minipage}%
%    \begin{minipage}[t]{.3\textwidth}
%      \begin{tikzpicture}[automaton]
%        \state[accepting, initial] (A) at (0,1) {$A$}; 
%        \state[accepting]          (B) at (1,1) {$B$}; 
%        \state[accepting]          (C) at (0,0) {$C$}; 
%        \state[accepting]          (D) at (1,0) {$D$}; 
% 
%        \path (D) edge[loop right] node {$a$}  (D);
%        \path (A) edge             node {$b$} (B);
%        \path (A) edge             node {$a$} (C);
%        \path (C) edge             node {$b$} (B);
%        \path (C) edge             node {$a$} (D);
%      \end{tikzpicture}
%    \end{minipage}
% 
%  \end{question}
% 
%\end{correction}
% 
%\begin{correction}
% 
%  \begin{question}
%  \item {\small
%    \begin{tabular}[t]{| l | c c c c |}
%      \hline
%      & A & B & C  & D \\
%      \hline
%      Accessible          & $\{1,2,3,4\}$ & $\{1, 2, 3, 4, 5, 7\}$  & $\{1, 2, 3, 4, 5, 6\}$  & $\{1, 2, 3, 4, 5, 6, 7, 8, 9, 10, 11\}$ \\
%      Inaccessible        & $\emptyset$   & $\{6\}$                 & $\emptyset$             & $\emptyset$                             \\
%      Co-accessible       & $\{1,2,3,4\}$ & $\{1, 2, 5\}$           & $\{1, 2, 3, 6\}$        & $\{1, 2, 3, 4, 5, 6, 7, 8, 9, 10, 11\}$ \\
%      Stérile             & $\emptyset$   & $\{3, 4, 6, 7\}$        & $\{4, 5\}$              & $\emptyset$                             \\
%      Utile               & $\{1,2,3,4\}$ & $\{1, 2, 5\}$           & $\{1, 2, 3, 6\}$        & $\{1, 2, 3, 4, 5, 6, 7, 8, 9, 10, 11\}$ \\
%      Inutile             & $\emptyset$   & $\{3, 4, 6, 7\}$        & $\{4, 5\}$              & $\emptyset$                             \\
%      Complet             & oui           & non                     & non                     & non                                     \\
%      Emondé              & oui           & non                     & non                     & oui                                     \\
%      Unitaire            & oui           & oui                     & oui                     & oui                                     \\
%      Standard            & non           & oui                     & non                     & oui                                     \\
%      Normalisé           & non           & non                     & non                     & non                                     \\
%      $\varepsilon$-libre & oui           & oui                     & oui                     & non                                     \\
%      Déterministe        & oui           & non                     & non                     & non                                     \\
%      \hline
%  \end{tabular}}
%  \item 
%    \begin{minipage}[t]{.4\textwidth}
%      $B$ sur l'alphabet $\{a,b,c,d\}$ :\\
% 
%      \scalebox{.9}{\begin{tikzpicture}[shorten >=1pt,node distance=1.7cm,on grid,auto]
%          \node[state, initial, initial text=, accepting] (1)              {$1$}; 
%          \node[state, accepting]                         (2) [right=of 1] {$2$}; 
%          \node[state, accepting]                         (5) [below=of 2] {$5$}; 
% 
%          \path[-latex] (2) edge[loop above, looseness=5] node {$a$}    (2);
%          \path[-latex] (5) edge[loop below, looseness=5] node {$b, c$}    (5);
%          \path[-latex] (1) edge node {$a$} (2);
%          \path[-latex] (1) edge node[swap] {$b, c$} (5);
%          \path[-latex] (2) edge[bend left=6mm] node {$b$} (5);
%          \path[-latex] (5) edge[bend left=6mm] node {$a$} (2);
%      \end{tikzpicture}}
%    \end{minipage}%
%    \begin{minipage}[t]{.4\textwidth}
%      $C$ sur l'alphabet $\{0,1,2,3,4,5,6,7,8,9\}$ :\\
%      \scalebox{.9}{\begin{tikzpicture}[shorten >=1pt,node distance=1.7cm,on grid,auto]
%          \node[state, initial, initial text=, accepting] (1)              {$1$}; 
%          \node[state]                                    (2) [right=of 1] {$2$}; 
%          \node[state]                                    (3) [right=of 2] {$3$}; 
%          \node[state, accepting]                         (6) [below=of 3] {$6$}; 
% 
%          \path[-latex] (1) edge[loop above, looseness=5] node {$0$}    (1);
%          \path[-latex] (1) edge node {$1$} (2);
%          \path[-latex] (2) edge node {$0$} (3);
%          \path[-latex] (3) edge[bend left=6mm] node {$0$} (6);
%          \path[-latex] (6) edge[bend left=6mm] node {$0$} (3);
%      \end{tikzpicture}}
%    \end{minipage}
%  \end{question}
% 
%\end{correction}
% 
%\endgroup
%\endinput
